\subsection{Stored bytes and allocated bytes}

Persistent row data are only part of device allocation: models, query
workspace and temporary copies are additional. For byte-sized primary
codes, row and list-offset storage are
\begin{equation}
 S_{\rm rows}=N\left(M+\left\lceil\frac{dB_r}{8}\right\rceil+8\right),
 \qquad S_{\rm offsets}=4(n_{\rm list}+1).
 \label{eq:4-row-storage}
\end{equation}
IDs and corrections contribute $8N$ bytes. With $M=d/8$, $B_r=8$,
the code payload alone is
\[
 b_{\rm code}=d/8+d=9d/8=1.125d.
\]
This exceeds 8-bit scalar quantisation (SQ8) by 12.5\%; graph storage
and vector-code bytes must therefore be considered together.

\begin{table}[!htbp]
\centering
\caption[Per-vector storage model]{Per-vector bytes. JHQ includes codes, ID and correction; SQ8 is
vector codes alone; CAGRA int8 models vector codes plus 32 four-byte
neighbours. Offsets, model parameters, workspace and other runtime
allocations are excluded. These components are not complete process totals.}
\label{tab:4-bytes}
\footnotesize\setlength{\tabcolsep}{4pt}
\begin{tabular}{@{}rrrrrrrr@{}}
\toprule
$d$ & Primary & Residual & ID + corr. & JHQ & SQ8 & \shortstack{CAGRA\\int8} & FP32 \\
\midrule
768 & 96 & 768 & 8 & 872 & 768 & 896 & 3,072 \\
1,024 & 128 & 1,024 & 8 & 1,160 & 1,024 & 1,152 & 4,096 \\
1,536 & 192 & 1,536 & 8 & 1,736 & 1,536 & 1,664 & 6,144 \\
3,072 & 384 & 3,072 & 8 & 3,464 & 3,072 & 3,200 & 12,288 \\
\bottomrule
\end{tabular}
\end{table}

At graph degree $d_g=32$, CAGRA int8 uses $d+4d_g$ row bytes. Stella's
JHQ row payload (19.20\,GiB) is close to int8 CAGRA's reported index
allocation of 19,656\,MiB, with different accounting boundaries.
Changing $\alpha$ changes residual reads, not the stored row bytes.

Basic model storage accounts for the rotation matrix, primary and
residual codebooks, and centroids with their norms:
\[
 S_{\rm model}
 =4d^2+4MKD_s+4MK_r+4n_{\rm list}(d+1)
 \quad\text{bytes}.
\]
Derived arrays are additional. Query workspace holds two query matrices,
centroid scores, primary LUTs, probes and retained candidates; refinement
scores residuals directly. The reported 128-probe configuration reserves
about 93.8\,MiB for scan groups even when unused
(\cref{eq:5-group-workspace}). Allocated capacity therefore also affects
resident memory.

The input remains caller-owned. For $n_b$ pipeline buffers with
$a_{\rm in}$ staged rows each, pinned host staging and device ordering
arrays use
\[
 S_{\rm pinned}=4n_ba_{\rm in}d,\qquad
 S_{\rm order}=S_{\rm inverse}=4N.
\]
The inverse map is temporary. Assignment scratch is freed before
gathering (\cref{eq:4-staging-rows,eq:4-assignment-scratch}), so host
staging, resident data and construction peaks have distinct lifetimes.

\subsection{Cache reuse and invalidation}

Cache keys omit code versions and full-data hashes, so code/data changes
require invalidation. Parameter sweeps reuse models; cold builds refit
them. Similar recall alone does not establish model identity. A cache
hit still incurs QR, model upload and add; warm-up and scheduler setup
are also needed before steady-state search. Construction and startup
timing boundaries are compared in \cref{sec:7-7}.

Changing $\alpha$ needs no retraining. Query buffers are reused by CUDA
Graph replay, and workspace grows before capture. With the result count
fixed, changes to batch size, survivors, probes, lanes, groups or bounds
trigger recapture. Calibration may bypass capture; index-owned buffers
remain allocated until release.
